\section*{Introduction}
Biomedical abstracts are a compact but information-dense substrate for evidence synthesis. In randomized-trial and clinical-research abstracts, sentences describing outcomes, findings, and conclusions are often the first targets for screening, triage, and downstream evidence extraction. A sentence-level model that can identify this rhetorical content is not a substitute for full critical appraisal, but it can reduce the cost of navigating large search results and make subsequent human review more efficient.

The methodological question is not only whether a model can obtain high discrimination. In healthcare-adjacent machine learning, calibrated probabilities and external validation are central safeguards against over-interpreting a model's apparent performance. A model that ranks sentences well may still provide misleading probability estimates, and a model tuned on one corpus may degrade when evaluated on another corpus with different labelling conventions. This study therefore treats calibration, external validation, and reproducibility as first-class outcomes.

CPU-feasible modelling is also scientifically important. Not every clinical informatics group or evidence-synthesis team has access to local GPU infrastructure, and not every biomedical NLP task requires a large Transformer model. Sparse linear models and lightweight recurrent neural networks remain attractive when a benchmark must be transparent, rerunnable, and inexpensive. This study asks whether such methods can solve a clinically meaningful abstract triage task under explicit CPU-only constraints.

The contributions are fourfold. First, the study defines a binary findings/conclusion detection task using public PubMed-PICO data and a label-compatible PubMed RCT external validation corpus. Second, it compares full-training TF-IDF baselines against CPU-limited neural models, including LSTM and BiLSTM variants. Third, it reports discrimination, thresholded classification performance, calibration metrics, reliability plots, and subgroup analysis from saved predictions. Fourth, it packages code, logs, metrics, tables, figures, and LaTeX source to support independent checking.
